\section{Proposed Methodology}
\label{sec:method}
\subsection{Context-gated GATv2--Temporal Transformer}
For each node and time step, every feature group $g$ is projected to a common width $d=64$,
$h_{t,v,g}=\mathrm{ReLU}(W_g x_{t,v,g})$. A scoring network gives $s_{t,v,g}=w^\top\tanh(U h_{t,v,g})$; scores of
unavailable groups ($m_{t,v,g}=0$) are masked to $-\infty$ and a softmax over the remaining groups yields weights
$\alpha_{t,v,g}$. The fused representation is $z_{t,v}=\sum_g\alpha_{t,v,g}h_{t,v,g}$; if no group is available a
learned ``unavailable'' vector is used and the event is logged. Two GATv2 layers~\cite{brody2022gatv2} (4 heads,
residual connections, layer normalisation) propagate $z$ over the fixed directed graph at every time step
(implemented with PyTorch Geometric~\cite{fey2019pyg}). Each node's sequence then receives a learned positional
embedding and passes through a 2-layer, 4-head Transformer encoder (feed-forward width 128, dropout 0.1); a linear head on
the last state outputs the three horizons. Targets are standardised per link with training-period statistics and the
masked Huber loss is averaged over observed targets only; batches without any observed target are skipped and
counted. Optimisation: AdamW (lr $10^{-3}$, weight decay $10^{-4}$), batch size 8, gradient clipping at 1.0, early
stopping on validation MAE with patience 5 and at most 30 epochs, seed 42. The temporal snapshots are held in a
PyTorch Geometric Temporal~\cite{rozemberczki2021pgt} \texttt{StaticGraphTemporalSignal}, from which training windows
are drawn lazily. The gate weights are not interpreted as causal importance.

\subsection{Baselines and ablations}
Persistence (last observation), a training-period historical average per link, day type and 5-minute slot, XGBoost and
LightGBM (one model per horizon on 12 speed lags and the same causal features), a GRU, and a temporal-only Transformer
(the proposed model without message passing). Ablations remove the gate (concatenation instead), or one feature group
at a time.

\subsection{Uncertainty and reliability gate}
With the selected model frozen, per-horizon absolute residuals on the calibration partition give the split-conformal
quantile $q_h$ as the $\lceil (n+1)(1-\alpha)\rceil$-th smallest residual ($\alpha=0.1$); a per-link variant is also
computed when at least 50 calibration residuals exist. The advisory gate defers when data are older than 20 minutes,
more than one third of recent traffic inputs are missing, the interval is wider than a threshold chosen on calibration
data (80th percentile of per-link widths), or calibration support is below 200 samples.

\subsection{Explanations}
For an alert (link $v$, horizon $h$, origin $t$) the explained scalar is the predicted speed in mph. Players are the
target's recent traffic (last 30 min), older traffic (30--120 min), the traffic of its graph neighbours and each context
group; an absent player takes values from 16 training-period background windows. All $2^P$ coalitions are evaluated
(exact KernelSHAP~\cite{lundberg2017shap}). We additionally measure neighbour-edge removal and temporal-window
sensitivities. Tree baselines are explained with TreeSHAP and labelled as baseline explanations.

\subsection{Retrieval, LLM and agentic workflow}
Approved documents (U.S. federal signal-timing guidance, SUMO documentation, and a clearly labelled project-defined
experimental policy) are hashed, allow-listed, chunked with a 200 word-piece target and indexed with a FAISS
inner-product index over normalised all-MiniLM-L6-v2 embeddings; chunks that still exceed the encoder's 256-token
window (171 of 1{,}633 in the built index, mostly tables of the PDF manual) are truncated by the encoder, and the
count is recorded in the index manifest. A query whose best passages
score below 0.35 returns INSUFFICIENT\_EVIDENCE. The LLM receives forecasts, intervals, quality flags, computed
attributions, retrieved passages (marked as untrusted text) and the allowed action catalogue, and must return JSON that
satisfies a Pydantic schema enforced through Ollama's structured-output option. A LangGraph state machine runs
quality $\rightarrow$ forecast $\rightarrow$ uncertainty $\rightarrow$ explanation $\rightarrow$ retrieval
$\rightarrow$ recommendation $\rightarrow$ safety $\rightarrow$ simulation $\rightarrow$ decision $\rightarrow$ human
review, with deterministic routing to DEFER\_DATA, DEFER\_UNCERTAINTY, DEFER\_EVIDENCE, REJECT\_SCHEMA,
REJECT\_CONSTRAINT, REJECT\_SIMULATION, PENDING\_HUMAN\_REVIEW or ERROR\_TOOL. Only the recommendation node calls the
LLM.

\subsection{Safety constraints and simulation gate}
Allowed actions are a bounded green extension (0--10 s) of an existing green phase, selection of a pre-validated
program, or monitoring only. Thirteen deterministic checks (mapping validity, signal and phase existence, action type,
bounds, minimum/maximum green, clearance preservation, pedestrian and side-road protection availability, data freshness,
forecast reliability, evidence support) plus a text policy filter must pass. Candidates are then applied through a
protective TraCI wrapper that refuses non-green phases and out-of-bound durations, and compared with the fixed-time
baseline on matched seeds; a candidate passes only if corridor delay improves by at least 1\%, side-road delay rises by
at most 10\% and no vehicle teleports. Each flow drives one fixed free-flow route so that demand stays on the corridor,
the fixed-time programs are re-timed for the routine demand with Webster's method (SUMO
\texttt{tlsCycleAdaptation.py}; a declared assumption, not NYC timing plans), and the testbed itself must first pass a
baseline test: the side-street demand level is the largest at which the
fixed-time baseline, on development seeds, has no teleport and at most 1\% unfinished trips; without such a level no
candidate is scored (NOT\_EVALUATED). Pedestrian performance is not modelled and is reported as NOT\_EVALUATED.
